\documentclass[11pt,a4paper,twocolumn]{article}

\usepackage[utf8]{inputenc}
\usepackage[T1]{fontenc}
\usepackage{lmodern}
\usepackage[a4paper,margin=1.9cm,columnsep=0.6cm]{geometry}
\usepackage{amsmath,amssymb,amsthm}
\usepackage{mathtools}
\usepackage{graphicx}
\usepackage{booktabs}
\usepackage{hyperref}
\usepackage{enumitem}
\usepackage{xcolor}

\hypersetup{colorlinks=true,linkcolor=blue,citecolor=blue,urlcolor=blue}

\newtheorem{theorem}{Theorem}
\newtheorem{lemma}[theorem]{Lemma}
\newtheorem{corollary}[theorem]{Corollary}
\newtheorem{proposition}[theorem]{Proposition}
\theoremstyle{definition}
\newtheorem{definition}[theorem]{Definition}
\newtheorem{axiom}[theorem]{Axiom}
\newtheorem{construction}[theorem]{Construction}
\theoremstyle{remark}
\newtheorem{remark}[theorem]{Remark}

\newcommand{\Feat}{\mathsf{Feat}}
\newcommand{\Capset}{\mathrm{Capset}}
\newcommand{\Src}{\mathcal{S}}
\newcommand{\Shape}{\mathrm{Shape}}
\newcommand{\Adm}{\mathrm{Adm}}
\newcommand{\Med}{\mathsf{m}}
\newcommand{\Gph}{\Gamma}
\newcommand{\V}{\mathcal{V}}
\newcommand{\E}{\mathcal{E}}
\newcommand{\wt}{w}
\newcommand{\sepc}{\sigma}
\newcommand{\flo}{\beta}
\newcommand{\Req}{\mathrm{Req}}
\newcommand{\Conf}{\mathrm{Conf}}
\newcommand{\LM}{\mathrm{LM}}

\title{Shape Is Not Admissibility:\\
A Capability-Theoretic Foundation for Data Modeling in Buhera OS}

\author{Kundai Farai Sachikonye\\
\small Technical University of Munich\\
\small AIMe Registry for Artificial Intelligence\\
\small \texttt{kundai.sachikonye@tum.de}}

\date{\today}

\begin{document}
\maketitle

\begin{abstract}
Data modeling languages---LinkML, JSON Schema, SHACL, OWL---answer the question
``what does a record of this kind look like.'' They do this well, and nothing in
this paper disputes it. We ask a different question: given a schema that answers
the first question completely and honestly, does an agent holding it know which
\emph{questions} are answerable against the data it describes? We show the answer
is no, in general, and that the gap is not a defect of any particular language but
a theorem about the relationship between two properties a data-modeling framework
can certify. We generalize the capability-declaration mechanism already proven
sound for chemical file-format conversion---a declared feature set $\Capset(\Src)
\subseteq \Feat$ with statically decidable containment---to data schemas in
general, and show that any shape language, LinkML included, instantiates this
mechanism faithfully: a schema \emph{is} a capability declaration, and nothing is
lost by treating it as one. We then prove the separation result that motivates
this paper: shape-certification and admissibility-certification are independent
properties, witnessed by two data sources with identical declared capability sets
and identical schema-conformant records, one of which admits a given question and
one of which does not, for reasons invisible to any capability set built from
per-field or per-class features alone. Admissibility is a property of a floor---a
global minimum separation cost over the data's contact structure---and no shape
language of the kind surveyed has a construct that names it. The consequence is
constructive, not critical: Buhera OS's data layer is the composition of (i) any
sufficient shape language, adopted or extended as the surface demands, satisfying
$\Feat$ and $\Capset$, and (ii) a floor/residue layer beneath it that certifies
admissibility directly. We give the composition law relating the two, validate it
against the existing ladder and representation-converter results, and close with
the precise sense in which adopting LinkML costs Buhera nothing and settles
nothing.
\end{abstract}

\tableofcontents

\section{Introduction}
\label{sec:intro}

A data-modeling language lets a community agree, once, on what a record of some
kind contains: a reaction has participants and a rate; an enzyme has a sequence
and an EC number; a dataset has a license and a provenance statement. LinkML
formalizes this agreement as a schema of classes, slots, types, and enumerations,
expressed in YAML, compiled by generators into JSON Schema, SHACL, OWL, SQL DDL,
and typed code~\cite{linkml}. NFDI4Cat's Chem-DCAT-AP is exactly such a schema,
extending DCAT-AP with PROV-O provenance obligations for catalysis datasets. None
of this is in question here, and none of it is what this paper is about.

What a schema gives a consuming agent is a certificate of \emph{shape}: this
record, if it validates, has the fields the schema declares, typed as the schema
declares them. What a schema does not give, and cannot be extended to give by
adding more classes, slots, or constraints, is a certificate of
\emph{admissibility}: whether a particular question posed against the data the
schema describes has an answer that the data can actually produce, as opposed to
one that is silently absent, structurally inexpressible in the schema's own
terms, or blocked by a resource that ran out before an answer was reached. A
companion result already proves this distinction for one concrete
case---SPARQL and SHACL over RDF triples cannot express admissibility, because
admissibility depends on a global minimum over the data's contact structure that
no locally-matchable pattern can see~\cite{ladder}. This paper generalizes that
result from one query language to the entire family of shape languages, and gives
Buhera OS the data-modeling primitive it has so far lacked: not a competing
schema language, but the missing complement to any schema language, stated as a
theorem about what shape can and cannot certify.

The paper's contributions:

\begin{itemize}[leftmargin=1.4em]
  \item A generalization of the capability-declaration mechanism proven for
  chemical format conversion~\cite{masamune}---a declared feature set
  $\Capset(\Src)\subseteq\Feat$ with statically decidable containment---from file
  formats to data schemas of any kind (Section~\ref{sec:capability}).
  \item A proof that LinkML schemas, and by the same construction JSON Schema,
  SHACL, and OWL class expressions, instantiate this mechanism faithfully: a
  schema's declared classes and slots \emph{are} a feature set and a capability
  declaration, not merely analogous to one (Section~\ref{sec:linkml}).
  \item The separation theorem: shape-certification and admissibility-
  certification are independent, witnessed by two capability-identical,
  schema-conformant data sources that disagree on the admissibility of the same
  question (Section~\ref{sec:separation}).
  \item A composition law giving Buhera's data layer as shape-plus-floor, with
  the precise conditions under which adopting an existing schema costs nothing
  and the precise conditions under which it is insufficient
  (Section~\ref{sec:composition}).
  \item Validation against the existing ladder and representation-converter
  implementations, which already exhibit this composition without having
  named it (Section~\ref{sec:validation}).
\end{itemize}

We say at the outset what this paper does not claim. It does not claim LinkML is
inadequate for what it is built to do, does not propose replacing it, and does
not propose a new syntax for classes and slots. Where a shape language is
sufficient for a surface Buhera must expose---and Chem-DCAT-AP interoperability
is exactly such a surface---the correct engineering decision is to adopt it
unmodified. The paper's entire content is the theorem that adoption of a shape
language, however complete, leaves a specific and namable question
unanswered, and the construction that answers it.

\section{The capability primitive, generalized}
\label{sec:capability}

\subsection{From file formats to data schemas}

The representation-converter work defines a \emph{source} as a grammar together
with a declared capability set drawn from a feature alphabet: $\Feat$ enumerates
structural features a format may state (element, connectivity, stereochemistry,
provenance, and so on for chemical structures), and $\Capset(\Src)\subseteq\Feat$
is what a given format's reader actually extracts~\cite{masamune}. Nothing in
that definition is specific to chemical file formats; it is specific to the
question ``what can this thing tell me,'' asked of any producer of structured
data.

\begin{definition}[Feature alphabet]
\label{def:feat}
A \emph{feature alphabet} $\Feat$ for a domain of discourse is a finite set of
named, independently checkable properties a record in that domain may or may not
state. Independence means membership of one feature in a capability set is
decided without reference to any other.
\end{definition}

\begin{definition}[Data source, capability set]
\label{def:source}
A \emph{data source} $\Src$ over $\Feat$ is a pair $(\Gamma_\Src, \Capset(\Src))$
where $\Gamma_\Src$ is a grammar or schema governing records the source can
produce, and $\Capset(\Src)\subseteq\Feat$ is the subset of features the source's
extraction procedure actually populates, as opposed to declares syntactically
available and leaves absent or defaulted.
\end{definition}

\begin{definition}[Request, containment]
\label{def:request}
A \emph{request} $\Req\subseteq\Feat$ is the feature set a consumer needs
populated to carry out some downstream operation. The request is
\emph{satisfiable} by $\Src$ iff $\Req\subseteq\Capset(\Src)$.
\end{definition}

\begin{theorem}[Static decidability of containment]
\label{thm:decidable}
Given $\Req$ and a declared $\Capset(\Src)$, satisfiability is decidable in
$O(|\Feat|)$ time, without reading any record the source produces.
\end{theorem}

\begin{proof}
Immediate set-containment check over two subsets of a finite, fixed alphabet;
this is Theorem~4.7 of~\cite{masamune}, restated for a general $\Feat$ rather
than the chemical-structure alphabet given there.
\end{proof}

\begin{remark}[What this buys, and what it costs]
Theorem~\ref{thm:decidable} is the entire engineering payoff of treating a
schema as a capability declaration: a consumer can refuse a request before
touching data, with a reason (the missing features), rather than after
discovering an empty or malformed result. The cost, stated precisely in
Section~\ref{sec:separation}, is that containment of a declared feature set says
nothing about a property that is not a feature of any single record---which is
exactly what admissibility turns out to be.
\end{remark}

\subsection{Honesty is assumed, not verified}

\begin{definition}[Honest declaration]
\label{def:honest}
A source's declaration $\Capset(\Src)$ is \emph{honest} if, for every feature
$f\in\Capset(\Src)$, the extraction procedure populates $f$ on every record for
which the underlying data determines a value, and \emph{never} populates $f$ with
a supplied or fabricated value tagged as though it were source-stated.
\end{definition}

\begin{proposition}[Honesty is not mechanically checkable]
\label{prop:honesty-unverifiable}
No procedure operating on $\Src$'s declared capability set alone decides whether
$\Src$ is honest. Honesty is checkable only by differential comparison against an
independently trusted extractor for the same underlying data.
\end{proposition}

\begin{proof}
This is Section~14 of~\cite{masamune}, restated: over-declaration is
indistinguishable from correct declaration at the interface level, since both
produce a populated field. We carry the assumption forward here as a standing
condition on every capability set discussed in this paper, exactly as
in~\cite{masamune}.
\end{proof}

\section{Shape languages instantiate the capability primitive}
\label{sec:linkml}

\subsection{LinkML as a capability declaration}

A LinkML schema declares a set of classes, each with a set of slots, each slot
typed and optionally constrained (cardinality, value range, enumeration
membership). We show this is, without approximation, an instance of
Definition~\ref{def:source}.

\begin{proposition}[LinkML schemas are capability declarations]
\label{prop:linkml-capset}
Let a LinkML schema $\LM$ declare classes $C_1,\dots,C_k$, each with slots
$\{s_{i1},\dots,s_{im_i}\}$. Define $\Feat_\LM = \{(C_i,s_{ij})\}$, the set of
qualified slot names. Then $\LM$, together with its generated validator, is a
data source $(\Gamma_\LM,\Capset(\LM))$ in the sense of
Definition~\ref{def:source}, with $\Capset(\LM)=\Feat_\LM$ exactly when every
declared slot is \texttt{required}, and $\Capset(\LM)\subsetneq\Feat_\LM$
whenever some slot is optional and, for a given instance record, absent.
\end{proposition}

\begin{proof}
Direct correspondence: $\Gamma_\LM$ is the schema's YAML grammar together with
its class/slot inheritance rules, which is exactly a grammar governing
producible records (Definition~\ref{def:source}). Required slots are populated
on every valid instance by LinkML's own validation semantics, hence lie in every
instance's realized capability set; optional slots may or may not be populated
per instance, giving the instance-dependent capability set claimed. Containment
checking (does a request for slots $\Req$ succeed against a given instance) is
exactly LinkML's JSON-Schema-generated validation, which is the
$O(|\Feat_\LM|)$ check of Theorem~\ref{thm:decidable} for this alphabet.
\end{proof}

\begin{corollary}[The generalization is uniform across the shape-language family]
\label{cor:family}
The same correspondence holds, mutatis mutandis, for a JSON Schema's declared
properties, a SHACL shape's declared property constraints, and an OWL class
expression's declared restrictions: in every case, $\Feat$ is the language's
finest declarable unit (property, slot, or restriction), and $\Capset$ is the
subset an instance actually satisfies.
\end{corollary}

\begin{proof}
Each of the three languages already compiles to or from LinkML by declared
generator correspondence~\cite{linkml}; composing Proposition~\ref{prop:linkml-capset}
with a generator's own round-trip guarantee transports the capability-set
correspondence unchanged.
\end{proof}

\begin{remark}[This is the whole of the ``LinkML question,'' honestly answered]
Proposition~\ref{prop:linkml-capset} and Corollary~\ref{cor:family} are the
complete, formal version of the informal claim that opened this line of work:
Buhera OS need not choose between adopting LinkML and building something of its
own, because LinkML already \emph{is} a (very good, tooled, community-adopted)
way of populating $\Feat$ and declaring $\Capset$ for a given surface. Adopting
Chem-DCAT-AP for a dataset-catalog surface is, under this correspondence,
adopting a specific, already-engineered $\Feat_\LM$---not a competing
foundational commitment. Nothing is given up by using it where it is sufficient.
\end{remark}

\section{Shape is not admissibility}
\label{sec:separation}

\subsection{What admissibility means}

We adopt admissibility from the ladder construction~\cite{ladder}, restated here
independently of catalysis.

\begin{definition}[Contact structure of a data source]
\label{def:contact-of-source}
A data source's records induce a contact graph $\Gph=(\V,\E,\wt,\Med)$: items
$\V$ are the source's addressable entities (rows, individuals, nodes---the
granularity is the source's own), edges $\E$ record a positive-weighted relation
of co-occurrence or reference between items, and $\Med$ is a distinguished medium
vertex representing the ambient population from which any item is drawn.
\end{definition}

\begin{definition}[Admissibility (restated)]
\label{def:admissibility}
A question ``is there a path of accountable relations from item $v_0$ to a
target class of item $x^\star$'' is \emph{admissible} at $v_0$ iff the separation
cost $\sepc(v_0,x^\star)=\min_{S\ni v_0,\,x^\star\notin S}\wt(\mathrm{cut}(S))$
does not exceed the source's floor $\flo$ by more than a stated accountability
margin $\varepsilon$, exactly as in Definition~6.2--6.3
of~\cite{ladder}.
\end{definition}

Admissibility is a property of a \emph{minimum over subsets} of the whole graph,
not a property attached to any single item, edge, or field.

\subsection{The separation theorem}

\begin{theorem}[Shape does not determine admissibility]
\label{thm:separation}
There exist two data sources $\Src_1,\Src_2$ over the same feature alphabet
$\Feat$, satisfying $\Capset(\Src_1)=\Capset(\Src_2)$, whose records validate
against the same shape schema (LinkML or otherwise) with no discernible
difference at the level of any single record, such that a question admissible
at some $v_0$ in $\Src_1$ is not admissible at the corresponding $v_0$ in
$\Src_2$.
\end{theorem}

\begin{proof}
This is exactly Theorem~6.4 of~\cite{ladder} (``Retrieval cannot express
admissibility''), applied here at one remove. The construction given there
exhibits two contact graphs with identical recorded triples---hence identical
per-record shape, since shape is read off individual records or their local
neighborhoods---in which a peripheral vertex's edge weight changes the
\emph{floor} $\beta$ without altering any pattern a shape schema, a SPARQL query,
or a SHACL shape can match. Since $\Feat$ and $\Capset$ (Definitions~\ref{def:feat}
and~\ref{def:source}) are themselves built from per-record, per-field
properties---this is what makes them independently checkable
(Definition~\ref{def:feat})---no assignment of $\Capset(\Src_1)=\Capset(\Src_2)$
can encode a global minimum. The two sources of~\cite{ladder} therefore also
satisfy $\Capset(\Src_1)=\Capset(\Src_2)$ under any feature alphabet built from
locally observable record properties, which includes every alphabet a class/slot
shape language of the kind in Section~\ref{sec:linkml} can express, by
Proposition~\ref{prop:linkml-capset}.
\end{proof}

\begin{corollary}[No enrichment of $\Feat$ repairs this]
\label{cor:no-enrichment}
Adding attributes, numeric constraints, or additional slots to $\Feat$---any
enrichment expressible as a per-record or per-local-neighborhood property---does
not close the gap of Theorem~\ref{thm:separation}.
\end{corollary}

\begin{proof}
This is Corollary~6.5 of~\cite{ladder} carried across the same correspondence
used in the proof of Theorem~\ref{thm:separation}: attributed triples and
arithmetic comparisons over them remain local, and $\min$ over an exponential
family of subsets has no representation as a conjunction of local
predicates~\cite{ladder}.
\end{proof}

\begin{remark}[Why this is not a complaint about LinkML]
Corollary~\ref{cor:no-enrichment} forecloses the natural rejoinder that a
sufficiently detailed schema would eventually capture admissibility. It cannot,
for a structural reason independent of how much care goes into the schema: the
floor is a global property and $\Feat$, by Definition~\ref{def:feat}, is built
from independently checkable---hence local---properties. This is a theorem about
what any shape language of this general kind can express, not an assessment of
how completely LinkML's authors have modeled catalysis metadata.
\end{remark}

\section{Composition: Buhera's data layer}
\label{sec:composition}

\subsection{The composition law}

\begin{definition}[Composed data layer]
\label{def:composed}
A \emph{composed data layer} for a surface is a pair $(\Src, \Gph_\Src)$: a
shape-language source $\Src=(\Gamma_\Src,\Capset(\Src))$ satisfying
Definition~\ref{def:source} (Section~\ref{sec:linkml} shows this may be an
adopted LinkML schema, extended or not), together with the contact graph
$\Gph_\Src$ of Definition~\ref{def:contact-of-source} induced by $\Src$'s
records.
\end{definition}

\begin{proposition}[The two certificates are independent and both necessary]
\label{prop:both-necessary}
A composed data layer $(\Src,\Gph_\Src)$ answers a request $\Req$ for shape by
Theorem~\ref{thm:decidable} applied to $\Capset(\Src)$, and answers a question
for admissibility by Definition~\ref{def:admissibility} applied to $\Gph_\Src$'s
floor. Neither computation is derivable from the other: Theorem~\ref{thm:separation}
forecloses deriving admissibility from $\Capset(\Src)$, and $\Capset(\Src)$ is
not recoverable from $\Gph_\Src$ alone, since the contact graph does not
retain the per-field typing distinctions the shape language draws (only the
existence and weight of a relation between items).
\end{proposition}

\begin{proof}
The first half is Theorem~\ref{thm:separation}. The second half follows because
Definition~\ref{def:contact-of-source} constructs $\Gph_\Src$ from items and
weighted relations only; two sources with different $\Feat$ and different
$\Capset(\Src)$ (different slot sets, different types) can induce
isomorphic contact graphs whenever the relations they record coincide, so
$\Gph_\Src$ alone underdetermines $\Capset(\Src)$.
\end{proof}

\begin{corollary}[Buhera's data-modeling primitive]
\label{cor:buhera-primitive}
Buhera OS's data-modeling layer is the composed data layer of
Definition~\ref{def:composed}, for whatever shape source a surface requires. Where
an existing shape language's $\Capset$ is sufficient for the surface's requests
(Chem-DCAT-AP for dataset-catalog interoperability, per~\cite{linkml}), it is
adopted unmodified; the floor/residue layer $\Gph_\Src$ is constructed
independently, on the same underlying items, and answers exactly the class of
question Section~\ref{sec:separation} shows the shape layer cannot.
\end{corollary}

\subsection{When adoption is sufficient, and when it is not}

\begin{proposition}[Sufficiency condition]
\label{prop:sufficiency}
Adopting an existing shape language $\LM$ unmodified is sufficient for a surface
iff every request the surface must serve is a shape request in the sense of
Definition~\ref{def:request}, and no consumer of the surface poses an
admissibility question in the sense of Definition~\ref{def:admissibility}.
\end{proposition}

\begin{proof}
Immediate from Proposition~\ref{prop:both-necessary}: shape requests are served
by $\Capset(\LM)$ alone; admissibility questions are not served by
$\Capset(\LM)$ at any degree of schema enrichment, by
Corollary~\ref{cor:no-enrichment}, and require $\Gph_\LM$.
\end{proof}

\begin{remark}[The practical reading, stated once and plainly]
A dataset-catalog listing, a metadata harvester, a FAIR-compliance check---these
are shape requests, and Chem-DCAT-AP already serves them well; extending or
modifying it for these purposes would add engineering cost for no theorem-backed
gain. ``Which enzymes have I not yet tried against this substrate, and would
trying them tell me anything the floor doesn't already foreclose''---the
catalytic-ladder questions already validated
in~\cite{ladder}---are admissibility questions, and no LinkML schema, however
extended, answers them. Buhera's contribution is not a rival to the first kind of
question; it is the only present construction that answers the second.
\end{remark}

\section{Validation}
\label{sec:validation}

We do not introduce a new validation suite here; instead we show the composition
of Section~\ref{sec:composition} is already realized, unnamed, by two existing
implementations, which this paper's constructions now let us describe precisely.

\begin{proposition}[The representation converter is a $\Capset$-only layer]
\label{val:masamune}
Masamune's verdict algebra (translated, incomplete, unsupported, malformed,
underdetermined, empty)~\cite{masamune} operates entirely at the level of
Definition~\ref{def:source} and Theorem~\ref{thm:decidable}: every verdict is
computed from $\Req$ and $\Capset(\Src)$ containment, before or independently of
any floor computation. It is, in the vocabulary of this paper, the shape half of
Definition~\ref{def:composed} alone, and its own stated limitation---honest
declaration is assumed, never verified (Proposition~\ref{prop:honesty-unverifiable})---
is exactly the limitation this paper inherits unchanged for the general case.
\end{proposition}

\begin{proposition}[The ladder architecture is a composed data layer]
\label{val:ladder}
The four-stage architecture of~\cite{ladder}---fetch via a plan language,
translate to a contact graph, apply the ladder, emit one of six verdicts
(answer, empty, unexpressed, unsupported, starved, exhausted)---is precisely
Definition~\ref{def:composed} realized for RDF/SPARQL sources: the plan
language's declared source capabilities are $\Capset(\Src)$
(Theorem~8.3 of~\cite{ladder}, restated here as
Theorem~\ref{thm:decidable}), and the ladder's rungs, floor, and role/direction
computations operate on $\Gph_\Src$. The \emph{unexpressed} verdict is the
architecture's own name for a request refused at the $\Capset$ stage; the
\emph{starved} verdict is a failure discovered only at the $\Gph_\Src$ stage.
Their coexistence as distinct, non-conflatable outcomes (Theorem~8.5 of
\cite{ladder}, non-degeneracy) is an existing, validated instance of
Proposition~\ref{prop:both-necessary}: shape and admissibility fail
independently and must be reported independently.
\end{proposition}

\begin{remark}[The live demonstration]
The federated-query front end already implementing this
distinction---rendering \emph{unexpressed}/\emph{unsupported} refusals with
zero requests issued as visually and semantically distinct from
\emph{starved} or \emph{empty} outcomes, against real biocatalysis questions---is
the demonstrable form of Corollary~\ref{cor:buhera-primitive}: a running system
in which shape-layer refusal and floor-layer refusal are already
kept apart, prior to this paper giving the distinction its general theorem.
\end{remark}

\section{Related work}
\label{sec:related}

LinkML~\cite{linkml} and the standards it compiles to and from---JSON Schema,
SHACL, OWL, DCAT-AP and its NFDI4Cat Chem-DCAT-AP
extension---are the shape languages this paper generalizes over
(Section~\ref{sec:linkml}); we take no position on their internal design and use
only the external interface (declare a shape, validate an instance against it)
that all of them share. The capability-set and verdict-algebra machinery is
Masamune's~\cite{masamune}, generalized here from chemical file formats to data
schemas at large. The admissibility construction, the floor, and the
separation result of Section~\ref{sec:separation} are the ladder
architecture's~\cite{ladder}, applied here one level of generality higher, from
SPARQL/SHACL specifically to the shape-language family as a whole. The
underlying contact-graph, floor, and residue mechanics are shared with the
trace calculus~\cite{trace} and ultimately derive from the occupation axiom
of~\cite{occupation}: a schema's declared shape is, in that vocabulary, a
catalogue of what can be individuated about a record; the floor is what remains
true of the record's population regardless of how finely any single record is
described.

\section{Conclusion}
\label{sec:conclusion}

A data-modeling framework answers ``what does a record look like,'' and Buhera OS
has not needed to build one of its own for that question, because good answers
already exist and are freely adoptable: LinkML's classes, slots, and generators
are, by Proposition~\ref{prop:linkml-capset}, already an instance of the
capability-declaration primitive this paper generalizes from chemical format
conversion. What Buhera OS lacked was not a competing answer to that question but
the theorem, given here as Theorem~\ref{thm:separation}, that a second question---
which of the questions I might pose against this data are actually
admissible---is answered by none of these languages at any degree of schema
enrichment, because the property in question is a global floor over the data's
contact structure and no shape language's declarable features are anything but
local. The corresponding construction, Corollary~\ref{cor:buhera-primitive}, is
Buhera's actual contribution to data modeling: not a fifth shape language beside
LinkML, JSON Schema, SHACL, and OWL, but the floor/residue layer that every one
of them is missing, composed underneath whichever of them is sufficient for a
given surface. Adopting LinkML for Chem-DCAT-AP interoperability, under this
account, costs nothing and settles nothing: it is the correct engineering choice
for the question it answers, and silent on the question this paper answers
instead.

\bibliographystyle{plain}
\bibliography{references}

\end{document}
